\documentclass[a4paper,10pt]{article}

\usepackage[T1]{fontenc}
\usepackage[utf8]{inputenc}

\usepackage{fontawesome}
\usepackage{latexsym}
\usepackage[empty]{fullpage}
\usepackage{titlesec}
\usepackage{marvosym}
\usepackage[usenames,dvipsnames]{color}
\usepackage{verbatim}
\usepackage{enumitem}
\usepackage[colorlinks,pdftex]{hyperref}
\usepackage{fancyhdr}
\usepackage{multicol}
\usepackage{eurosym}
\usepackage{xcolor}
\usepackage{xfrac}
\usepackage{vwcol} 
\usepackage{adjustbox}
\usepackage{longtable}
\usepackage{array}
\usepackage[french]{babel}
\definecolor{myBlue}{rgb}{0,0,0.55}


\hypersetup{
    urlcolor = {[rgb]{0,0,0}}
}

\newcommand\dunderline[3][-1pt]{{\%
  \sbox0{#3}% 
  \ooalign{\copy0\cr\rule[\dimexpr#1-#2\relax]{\wd0}{#2}}}}

\pagestyle{fancy}
\fancyhf{} 
\fancyfoot{}
\renewcommand{\headrulewidth}{0pt}
\renewcommand{\footrulewidth}{0pt}

\usepackage{geometry}
\geometry{top=0.2cm, bottom=0.4cm, left=0.3cm, right=0.3cm}

\urlstyle{same}

\raggedbottom
\raggedright
\setlength{\tabcolsep}{0in}

% Sections formatting
\titleformat{\section}{
  \vspace{-1pt}\scshape\raggedright\large
}{}{0em}{}[\[\color{black}\titlerule \vspace{-5pt}]

%-------------------------
% Custom commands
\newcommand{\resumeItem}[2]{
  \item\small{
    \textbf{#1}{: #2\vspace{-0.8pt}}
  }
}

\newcommand{\resumeSubheading}[4]{
  \vspace{1.15pt}
  \small
  \item
    \begin{tabular*}{0.97\textwidth}{l @{\extracolsep{\fill}}r}
      \normalsize{\textbf{#1}}   & #2 \\
      \small{#3} & #4 \\
    \end{tabular*}
  \vspace{1.5pt}
}

\def\Plus{\texttt{+\space}}

% -------- ORIGINAL ------------

\newcommand{\resumeSubItem}[2]{
    \small{\resumeItem{#1}{#2}\vspace{-4pt}}
	
}


\newcommand{\resumeSubHeadingListStart}{\begin{itemize}[leftmargin=*]}
\newcommand{\resumeSubHeadingListEnd}{\end{itemize}}
\newcommand{\resumeItemListStart}{\begin{itemize}}
\newcommand{\resumeItemListEnd}{\end{itemize}\vspace{-2pt}}


\newcommand{\resumeProjectListStart}{\begin{minipage}{0.97\textwidth} 
                                     \begin{itemize}[leftmargin=5mm]
                                     \vspace{3pt}}
\newcommand{\resumeProjectListEnd}{\end{itemize} 
                                   \end{minipage}}
\newcommand{\resumeProjectItem}[1]{
    \item\small{#1}
    }

% New commands for detailed version
\newcommand{\resumeDetailedProject}[5]{
    \item\small{
        \textbf{#1} \hfill \textit{#2} \\ 
        \begin{minipage}{0.95\textwidth}
            \small{\textbf{Langages :} #3 \\ 
            \textbf{Outils/Packages :} #4 \\ 
            \textbf{Difficult\'es :} #5}
        \end{minipage}
    }
    \vspace{3pt}
}

%-------------------------------------------
%%%%%%  CV STARTS HERE  %%%%%%%%%%%%%%%%%%%%%%%%%%%%


\begin{document}


%----------HEADING-----------------
\begin{center}
  \begin{minipage}{0.97\textwidth}
      \begin{minipage}[c]{0.6\linewidth}
          \huge \textbf{Ulysse \textsc{Boucherie}}
          \\
          \emph{\large{Ing\'enieur DevOps \& Data Scientist (MSc.)}}
      \end{minipage}
      \hfill
      \begin{minipage}[c]{0.215\linewidth}
          \begin{tabular*}{\textwidth}{c@{\extracolsep{\fill}} l}
              \href{mailto:uboucherie1@gmail.com}{\raisebox{-0.055\height}{\faEnvelope}} & \href{mailto:uboucherie1@gmail.com}{\small{uboucherie1@gmail.com}} \\ 
              \href{https://github.com/odysseu}{\raisebox{-0.035\height}{\faGlobe}} & \href{https://odysseu.github.io/}{\small{odysseu.github.io}} \\ 
              \href{https://github.com/odysseu}{\raisebox{-0.035\height}{\faGithub}} & \href{https://github.com/odysseu}{\small{odysseu}} \\ 
              \href{https://www.linkedin.com/in/ulysse-boucherie-7a3b2b111/}{\raisebox{-0.0925\height}{\faLinkedinSquare}} & \href{https://www.linkedin.com/in/ulysse-boucherie-7a3b2b111/}{\small{in/ulysse.boucherie}} \\ 
              {\raisebox{-0.0925\height}{\faHome}} & \small{Fontainebleau, France}
          \end{tabular*}
      \end{minipage}
  \end{minipage}
  \end{center}


%------------PROJET-----------------------

\begin{center}
\begin{minipage}{0.5\textwidth}
    \centering
    \large \textbf{Ing\'enieur DevOps \& data au service de la modernisation des Services Statistiques Publics}
\end{minipage}
\end{center}

\vspace{-5mm}

%-----------EXPERIENCE-----------------
\section{Experience}

\resumeSubHeadingListStart

\vspace{0.75mm}


\resumeSubheading{Minist\`ere de la s\'ecurit\'e int\'erieur}{Paris, France}{Data Scientist et r\'ef\'erent technique informatique}{2024 - aujourd'hui}

Aide \`a la Data Engineering, traitement, fiabilisation et automatisation de donn\'ees. Exploitation de donn\'ees sensibles \`a grande \'echelle en environnement s\'ecuris\'e (RGPD). SQL avanc\'e et optimis\'e (PostgreSQL, DuckDB), Python (pandas, numpy, pyarrow, SQLAlchemy). Construction et am\'elioration de flux de donn\'ees (ETL) avec Apache Airflow. Gouvernance et contr\^ole qualit\'e des donn\'ees. Production d'indicateurs d'aide \`a la d\'ecision publique.

Responsable des publications sur la d\'elinquance \'economique et financi\`ere
\resumeProjectListStart

\resumeDetailedProject
    {Refonte des processus ETL}
    {2024-Pr\'esent}
    {Python, SQL, R}
    {Apache Airflow, Pandas, NumPy, DuckDB, PostgreSQL}
    {Syst\`emes h\'erit\'es avec 20+ ans de dette technique ; R\'esistance au changement de la part des parties prenantes non techniques ; Probl\`emes de qualit\'e des donn\'ees dans les syst\`emes sources ; N\'ecessit\'e de maintenir la production pendant la migration}
    {Refonte des processus ETL pour 20\Plus agents, permettant de r\'eduire de plus de 80\% le d\'elai entre les acquisitions de donn\'ees et la mise en production. Mise en place de traitement incr\'ementiel, de cadres de validation des donn\'ees et d'alertes automatis\'ees. Formation des membres de l'\'equipe aux nouveaux processus et outils.}

\resumeDetailedProject
    {Mise en place d'une culture du changement}
    {2024-Pr\'esent}
    {Python, Markdown, Bash}
    {Git, GitHub, Docker, Jupyter, VS Code}
    {Public h\'et\'erog\`ene avec des comp\'etences techniques vari\'ees (de Excel uniquement \`a Python avanc\'e) ; r\'esistance \`a l'adoption de nouveaux outils ; besoin de d\'emontrer une valeur imm\'ediate}
    {Mise en place d'une culture du changement avec des outils et m\'ethodes de travail plus modernes. Cr\'eation de documentation, de mod\`eles et de kits de d\'emarrage. Organisation de sessions de formation et d'ateliers. D\'eveloppement d'outils internes pour r\'eduire la barri\`ere \`a l'entr\'ee pour les utilisateurs non techniques.}

\resumeDetailedProject
    {Outils open-source de publication}
    {2024-Pr\'esent}
    {Python, TypeScript, JavaScript}
    {React, FastAPI, Pandas, GitHub Actions, Docker}
    {\'Equilibre entre flexibilit\'e et standardisation ; garantie de la reproductibilit\'e ; gestion des d\'ependances ; consid\'erations de s\'ecurit\'e pour les outils destin\'es au grand public}
    {D\'eveloppement d'outils internes puis open source pour la publication industrialis\'ee d'articles et de contenus ouverts au grand public, r\'eduisant la dette technique et la d\'ependance. Mise en place de pipelines CI/CD, de tests automatis\'es et de g\'en\'eration de documentation.}

\resumeDetailedProject
    {Mod\`eles LLM pour le clustering de texte}
    {2024-Pr\'esent}
    {Python}
    {Transformers, scikit-learn, Sentence-BERT, HuggingFace}
    {Contraintes de confidentialit\'e (NDA) ; traitement de texte fran\c{c}ais avec jargon technique ; s\'election et fine-tuning de mod\`eles ; optimisation des performances}
    {Co-design, d\'eveloppement et d\'eploiement de mod\`eles LLM pour clusteriser les r\'eponses d'enqu\^ete \`a champs libres. Mise en \oe uvre de l'ing\'enierie de prompts, \'evaluation de plusieurs architectures de mod\`eles et d\'eploiement de la solution en environnement de production.}

\resumeDetailedProject
    {\'Etudes sur la d\'elinquance \'economique et financi\`ere}
    {2024-Pr\'esent}
    {Python, R, SQL}
    {Pandas, ggplot2, dplyr, PostgreSQL, DuckDB}
    {Donn\'ees provenant de multiples sources (Police Nationale, Gendarmerie) avec diff\'erents formats ; traitement de donn\'ees sensibles ; besoin d'analyses reproductibles}
    {R\'edaction d'\'etudes sur la d\'elinquance \'economique et financi\`ere en utilisant les donn\'ees de la Police Nationale et de la Gendarmerie. D\'eveloppement de pipelines d'analyse standardis\'es et de mod\`eles de visualisation.}

\resumeProjectListEnd

\resumeSubheading{INSEE}{Nantes, France}{Ing\'enieur DevOps}{2021 -- 2024}

Support et exploitation de syst\`emes d'information critiques. POC avec technologies avanc\'ees de Python (pandas, numpy, Dask, scikit-learn, XGBoost, TensorFlow/PyTorch, MLflow, Polars, pyarrow, Modin, Ray, FastAPI). Support et maintien de syst\`emes d'information critiques : stack compl\`ete en environnement conteneuris\'e (Docker, Kubernetes, Helm) avec couches basses (Linux Ubuntu/Debian, r\'eseaux TCP/IP, stockage S3) et hautes (PostgreSQL, Redis, microservices). Analyse et r\'esolution d'incidents applicatifs internes et open source. Am\'elioration continue des flux et processus SI avec observabilit\'e (Prometheus, Grafana, ELK, OpenTelemetry).

DevOps et support de la plateforme de datascience de l'INSEE
\resumeProjectListStart

\resumeDetailedProject
    {D\'eveloppement de la plateforme Onyxia}
    {2021-2024}
    {Go, Python, TypeScript, YAML}
    {Kubernetes, Helm, Docker, GitLab CI/CD, GitHub Actions, ArgoCD, Vault, Prometheus, Grafana}
    {Architecture multi-locataires complexe ; exigences de s\'ecurit\'e pour les donn\'ees gouvernementales ; int\'egration avec les syst\`emes h\'erit\'es ; besoins de scalabilit\'e}
    {Co-d\'eveloppement et installation d'une \textit{Plateforme As A Service} (\href{https://github.com/InseeFrLab/onyxia}{\textit{Onyxia}}). Conception d'architecture microservices, mise en \oe uvre de l'authentification/autorisation, et d\'eploiement de la stack de monitoring.}

\resumeDetailedProject
    {Impl\'ementation d'une stack technique moderne}
    {2021-2024}
    {Bash, Python, YAML}
    {Docker, Helm, Pipelines GitLab, GitHub Actions, Grafana, Prometheus, Vault, Terraform, Ubuntu/Debian, S3}
    {R\'esistance au changement de la part de l'\'equipe existante ; n\'ecessit\'e de maintenir les anciens syst\`emes pendant la transition ; besoins de formation pour 500+ utilisateurs}
    {D\'eveloppement et remplacement de la stack technique compl\`ete et neuve. Cr\'eation de guides de migration, automatisation des processus de d\'eploiement et r\'eduction du temps de d\'eploiement de plusieurs semaines \`a quelques heures (acc\'el\'eration par 10).}

\resumeDetailedProject
    {Support aux Data Scientists}
    {2021-2024}
    {Python, R, Bash}
    {JupyterHub, RStudio Server, PyCharm, VS Code, Kubernetes}
    {Large \'eventail de niveaux de comp\'etence des utilisateurs (d\'ebutants \`a avanc\'es) ; cas d'usage vari\'es (notebooks, tableaux de bord, IDE) ; contraintes de ressources}
    {Support de proximit\'e aux \textit{Data Scientists} d\'ebutants jusqu'aux avanc\'es pour l'int\'egration de leurs services (notebooks, applications internes de dataviz, IDE data : Rstudio \& Pycharm). D\'eveloppement de mod\`eles, de documentation et de bonnes pratiques.}

\resumeProjectListEnd

\resumeSubheading{INSEE}{Nantes, France}{Ing\'enieur data plateforme}{2020 -- 2021}

Sys-admin de la plateforme de datascience interne
\resumeProjectListStart

\resumeDetailedProject
    {D\'eploiement de composants de plateforme}
    {2020-2021}
    {PowerShell, Batch, Python}
    {Windows Server, Active Directory, IIS, SQL Server}
    {Infrastructure h\'erit\'ee avec des ressources limit\'ees ; contraintes de s\'ecurit\'e ; besoin de haute disponibilit\'e ; gestion complexe des utilisateurs}
    {Mise en place de briques permettant \`a 500\Plus datascientists de l'INSEE de gagner jusqu'\`a 80\% de leur temps pour d\'eployer des outils et \'etudes internes. Automatisation des scripts de d\'eploiement et cr\'eation de portails en libre-service.}

\resumeDetailedProject
    {Syst\`eme de monitoring des ressources}
    {2020-2021}
    {Python, PowerShell}
    {Prometheus, Grafana, Compteurs de performance Windows}
    {Visibilit\'e limit\'ee sur l'utilisation des ressources ; ressources physiques partag\'ees ; priorit\'es concurrentes ; besoin de donn\'ees historiques}
    {Mise en place de monitoring et de supervision des processus pour cadencer les ressources physiques partag\'ees (windows server). Optimisation de 20\% des resources utilis\'ees. Mise en place d'outils d'alerte et de planification de la capacit\'e.}

\resumeDetailedProject
    {Initiative de traduction de code}
    {2020-2021}
    {Python, R}
    {Pandas, tidyverse, reticulate}
    {Code propri\'etaire de 20 ans ; documentation limit\'ee ; r\'esistance des utilisateurs \`a l'aise avec l'ancien syst\`eme ; exigences de compatibilit\'e}
    {Mise en place de projets de traduction de code pour sortir d'une licence propri\'etaire de data-analyse implant\'ee depuis 20 ans. D\'eveloppement d'outils de migration et formation aux alternatives open-source.}

\resumeProjectListEnd

\resumeSubheading{Central Statistics Office}{Dublin, Irelande}{Stage de fin d'\'etudes}{2020}
Assistant Statisticien
\resumeProjectListStart

\resumeDetailedProject
    {Rapport officiel des statistiques annuelles}
    {2020}
    {R, LaTeX}
    {knitr, ggplot2, dplyr, R Markdown}
    {D\'elai serr\'e ; besoin de mise en forme officielle du gouvernement ; collaboration avec plusieurs parties prenantes ; exigences de validation des donn\'ees}
    {Co-\'edition du rapport officiel sur les statistiques annuelles de l'Irlande. D\'eveloppement de pipeline de g\'en\'eration de rapports automatis\'es et cr\'eation de mod\`eles r\'eutilisables pour les rapports futurs.}

\resumeDetailedProject
    {Traitement des donn\'ees m\'et\'eorologiques historiques}
    {2020}
    {R, Python}
    {Pandas, expressions r\'eguli\`eres, OCR (Tesseract)}
    {Enregistrements manuscrits de 250+ ans ; formats de donn\'ees incoh\'erents ; d\'efis de reconnaissance de caractères ; formats de dates historiques}
    {Mise en place d'un cahier des charges pour la lecture automatique de relev\'es m\'et\'eorologiques mensuels depuis 250\Plus ans. Impl\'ementation de pipeline OCR prototype et de r\`egles de validation des donn\'ees.}

\resumeProjectListEnd

\resumeSubHeadingListEnd


\vspace{-5mm}

%-----------EDUCATION-----------------
\section{Education}
\resumeSubHeadingListStart
\vspace{0.75mm}


\resumeSubheading{ENSAI}{Rennes, France}{Engineering Degree (MSc.)}{2018 -- 2021}
\resumeProjectListStart
\resumeProjectItem{Sp\'ecialit\'e informatique - Big Data - optimisation - \'economie (micro \& macro) - D\'el\'egu\'e de classe}
\resumeProjectListEnd

\resumeSubheading{Lyc\'ee Fran\c{c}ais Premier}{Fontainebleau, France}{Classe Pr\'eparatoire aux Grandes \'Ecoles}{2015 -- 2018}
\resumeProjectListStart
\resumeProjectItem{\textit{MPSI -- MP}: Math\'ematiques - Physique - Science de l'ing\'enieur - Sp\'ecialit\'e informatique - D\'el\'egu\'e de classe.}
\resumeProjectListEnd

\resumeSubHeadingListEnd


\vspace{-4mm}

%-----------COMPETENCES TECHNIQUES + LANGUES -----------------
\begin{minipage}[t]{0.65\linewidth}
    \section{Comp\'etences techniques}
    \resumeSubHeadingListStart
    \resumeSubItem{Programmation}{Python (avanc\'e), Java, JavaScript/TypeScript, R, Bash, Go}
    \resumeSubItem{Infrastructure}{Docker, Kubernetes, Helm, CI/CD (GitLab, GitHub Actions), Terraform}
    \resumeSubItem{Bases de donn\'ees}{PostgreSQL, DuckDB, MySQL}
    \resumeSubItem{Stockage}{S3}
    \resumeSubItem{Syst\`emes}{Linux (Ubuntu/Debian), Windows Server}
    \resumeSubItem{Plateformes Cloud}{Google Cloud Platform (GCP), Amazon Web Services (AWS)}
    \resumeSubItem{Ing\'enierie des donn\'ees}{Apache Airflow, Pandas, NumPy, Apache Spark, Dask, Polars, pyarrow, Modin, Ray}
    \resumeSubItem{ML/IA}{scikit-learn, TensorFlow, Transformers, HuggingFace, Fine-tuning de LLM, MLflow}
    \resumeSubItem{Monitoring}{Prometheus, Grafana, Stack ELK, OpenTelemetry, Datadog}
    \resumeSubItem{Outils DevOps}{Ansible, Vault, ArgoCD, GitHub Workflows}
    \resumeSubHeadingListEnd
    \end{minipage}
\hfill
\begin{minipage}[t]{0.3\linewidth}
    \section{Langues}
    \resumeSubHeadingListStart
    \resumeSubItem{Fran\c{c}ais}{Natif}
    \resumeSubItem{Anglais}{Bilingue}
    \resumeSubItem{Espagnol}{Courant}
    \resumeSubHeadingListEnd
\end{minipage}


%\vspace{0.5mm}

%-----------PROJECTS-----------------
\section{Projets personnels et int\'er\^ets}
\vspace{0.5mm}
\begin{minipage}[t]{0.9865\textwidth}
\resumeSubHeadingListStart

\resumeDetailedProject
    {HabitatCalc}
    {2022-Pr\'esent}
    {JavaScript, TypeScript, HTML/CSS}
    {React, Chart.js, Bootstrap, GitHub Pages}
    {Calculs financiers complexes ; gestion des cas particuliers ; design responsive pour mobile ; validation des entr\'ees utilisateur}
    {Application immobili\`ere : Optimisation entre location et achat d'un bien immobilier avec des param\`etres r\'eels. \href{https://odysseu.github.io/HabitatCalc/}{\textit{HabitatCalc}}.} 

\resumeDetailedProject
    {TrancheMaster}
    {2022-Pr\'esent}
    {JavaScript, TypeScript, HTML/CSS}
    {React, D3.js, GitHub Pages}
    {Calculs complexes de tranches d'imposition ; gestion des taux progressifs ; d\'efis de visualisation ; cas particuliers de la loi fiscale}
    {Application fiscale : Compr\'ehension des tranches marginales d'imposition sur le revenu et le taux d'imposition. \href{https://odysseu.github.io/TrancheMaster/}{\textit{TrancheMaster}}.} 

  \resumeSubItem{Veille et s\'eminaires}{\href{https://www.numerique.gouv.fr/agenda/etat-dans-le-nuage-2026/}{\textit{L'Etat dans le nuage 2026}}, Webinaires INSEE et Services Statistique Publique 2024 \`a 2026, Salon de l'IA 2023, KubeCon 2023, DevFest 2022 \& 2023, VivaTech 2022}
  \resumeSubItem{Sport}{Triathlon, volley-ball, ski alpin, escalade de blocs}
  \resumeSubItem{Musique}{piano (4 ans), batterie (10 ans)}
\resumeSubHeadingListEnd
\end{minipage}


%-------------------------------------------
\end{document}
